\chapter*{Introduction Générale}
\markboth{Introduction Générale}{Introduction Générale}
\addcontentsline{toc}{chapter}{Introduction Générale}

L’évolution rapide des technologies numériques et l’augmentation exponentielle du volume des données générées par les institutions financières ont profondément transformé les méthodes de gestion et de pilotage stratégique au sein du secteur bancaire. Dans ce contexte d'hyper-digitalisation, la donnée s'est imposée comme le socle fondamental sur lequel reposent l'évaluation des risques, la personnalisation des services clients, et la conformité aux exigences réglementaires de plus en plus strictes (telles que la norme BCBS 239, le RGPD ou encore la loi 09-08 de la CNDP au Maroc). Toutefois, malgré la mise en place d'infrastructures Big Data complexes telles que les Data Lakes pour centraliser ces informations, la seule capacité de stockage et d'ingestion ne suffit plus. Les données brutes restent souvent inexploitables ou dangereuses pour la prise de décision si elles s'avèrent incomplètes, incohérentes ou obsolètes.

Parallèlement, les progrès réalisés dans le domaine de l'ingénierie des données et de l'automatisation logicielle ouvrent de nouvelles perspectives pour assainir ces écosystèmes. Des outils modernes de validation des données, couplés à la puissance de calcul distribué (comme Apache Spark) et à des mécanismes d'orchestration avancés, permettent désormais d'automatiser l'audit continu des flux d'information. Ces technologies permettent non seulement de détecter les anomalies dès la source, d'enrichir les métadonnées de qualité, mais aussi de fournir des indicateurs de performance fiables (KPIs) avant que les données erronées n'impactent les modèles analytiques ou les reportings réglementaires en aval.

C’est dans cette optique que s’inscrit ce projet de fin d’études réalisé au sein du Data Office de CIH BANK. Son objectif consiste à concevoir, développer et industrialiser un pipeline automatisé de contrôle de la qualité des données (Data Quality) intégré directement à l'infrastructure Hadoop de la banque. En ciblant prioritairement le domaine stratégique des Tiers (incluant les informations d'identification KYC et les catégories socio-professionnelles CSP), le but est d'automatiser la validation, la surveillance et le monitoring de la qualité des données à travers des règles métiers strictes.

La solution développée repose sur une architecture Big Data complète comprenant l'ingestion massive des données depuis les systèmes opérants via Apache Sqoop, l'analyse et la validation des contenus à l'aide du framework Soda Core intégré à Apache Spark, ainsi que le stockage structuré des résultats dans des tables Hive. Au-delà du contrôle fonctionnel, ce projet met également l'accent sur la virtualisation des données via Dremio et la restitution interactive des métriques de qualité à travers un tableau de bord développé sous Tableau Desktop. Une chaîne d'orchestration complète basée sur Apache Airflow a été mise en œuvre afin d'automatiser ce cycle d'audit quotidiennement, garantissant ainsi une détection proactive des anomalies, une réduction des interventions manuelles chronophages, et une fiabilité accrue lors de l'exploitation des données par les métiers.

Le présent mémoire est organisé en quatre chapitres. Le premier chapitre présente le contexte général du projet, la présentation de l'organisme d'accueil CIH BANK, l'étude de l'architecture Big Data existante ainsi que la méthodologie de travail adoptée. Le deuxième chapitre est consacré à l’analyse détaillée des besoins fonctionnels et non fonctionnels, suivie de la conception des processus et règles de qualité selon les dimensions standards. Le troisième chapitre décrit l'architecture technique, les choix technologiques retenus (Hadoop, Spark, Soda Core, Airflow, Dremio, Tableau) ainsi que les modèles de données de restitution. Enfin, le quatrième chapitre détaille la réalisation technique, l’industrialisation du pipeline et la conception des tableaux de bord, avant de conclure sur le bilan du projet et ses perspectives d’évolution.
